% © 2025 Ing. David Jaroš — CC BY-NC-ND 4.0
%
% This work is licensed under a Creative Commons Attribution-NonCommercial-NoDerivatives
% 4.0 International License (CC BY-NC-ND 4.0).
%
% License History: Earlier drafts (up to v0.3) were released under CC BY 4.0.
% From v0.4 onward, all material is released under CC BY-NC-ND 4.0 to protect
% the integrity of the theoretical work during ongoing academic development.

\documentclass[11pt,a4paper]{article}
\usepackage{amsmath, amssymb, amsthm}
\usepackage{hyperref}
\usepackage{geometry}
\geometry{margin=2.5cm}

\newtheorem{theorem}{Theorem}
\newtheorem{lemma}{Lemma}
\newtheorem{proposition}{Proposition}
\newtheorem{definition}{Definition}
\newtheorem{remark}{Remark}

% Proof-level macros
\newcommand{\Lzero}{\textbf{[L0]}}
\newcommand{\Lone}{\textbf{[L1]}}
\newcommand{\Ltwo}{\textbf{[L2]}}
\newcommand{\SE}{\textbf{[SE]}}
\newcommand{\OPEN}{\textbf{[OPEN]}}
\newcommand{\PROVED}{\textbf{PROVED}}

\title{QED Complete Chain\\
\large Unified Biquaternion Theory — Canonical Bridges Series}
\author{Ing.\ David Jaroš}
\date{2025}

\begin{document}
\maketitle

\begin{abstract}
This document presents the complete derivation chain from the UBT field structure
to Quantum Electrodynamics (QED) in one continuous document.
Starting from the biquaternionic field $\Theta(q,\tau)$, we derive in sequence:
the U(1) gauge symmetry, the electromagnetic field tensor, Maxwell's equations,
the fermion coupling, and the Dirac equation limit.
Each step references its primary source and carries an explicit proof-level label.
This document is the single entry point for an external reviewer who wishes to
assess the QED sector of the Unified Biquaternion Theory.

\medskip
\noindent\textbf{Proof-level labels}: \Lzero{} exact algebraic identity;
\Lone{} proved given verified assumptions; \Ltwo{} requires additional lemmas;
\SE{} semi-empirical; \OPEN{} open problem.
\end{abstract}

\tableofcontents
\newpage

% ------------------------------------------------------------------
\section{Overview of the Derivation Chain}
% ------------------------------------------------------------------

The chain has five steps:

\begin{enumerate}
  \item \textbf{U(1) gauge symmetry} emerges from the scalar phase degree of
        freedom of $\Theta$.
  \item \textbf{Electromagnetic field tensor} $F_{\mu\nu}$ is defined as the
        curvature of the U(1) connection.
  \item \textbf{Maxwell equations} are the Euler--Lagrange equations of the
        QED Lagrangian in the classical limit.
  \item \textbf{Fermion coupling} enters through the gauge-covariant derivative
        acting on spinor excitations of $\Theta$.
  \item \textbf{Dirac equation} is the equation of motion for spinor modes in
        the U(1) background.
\end{enumerate}

\noindent\textbf{Status summary (before reading the full chain):}
\begin{center}
\begin{tabular}{lll}
\hline
\textbf{Step} & \textbf{Claim} & \textbf{Status} \\
\hline
1 & U(1) from right phase of $\Theta$ & \PROVED{} \Lzero \\
2 & $F_{\mu\nu}$ from U(1) connection & \PROVED{} \Lzero \\
3 & Maxwell equations from QED Lagrangian & \PROVED{} \Lone \\
4 & Minimal fermion coupling via $D_\mu$ & \PROVED{} \Lone \\
5 & Dirac equation as spinor EOM & \PROVED{} \Lone \\
\hline
\end{tabular}
\end{center}

% ------------------------------------------------------------------
\section{Step 1: U(1) Gauge Symmetry from the Biquaternionic Field}
% ------------------------------------------------------------------

\noindent\textbf{Status: \PROVED{} \Lzero}

\noindent\textbf{Primary sources:}
\texttt{canonical/interactions/sm\_gauge.tex}, Theorem G.3;
\texttt{canonical/fields/biquaternion\_algebra.tex}, \S2.

\medskip

\begin{definition}[Biquaternionic field]
The fundamental dynamical object of UBT is the field
\begin{equation}
  \Theta : \mathbb{R}^{1,3}\times\mathbb{C} \;\longrightarrow\;
  \mathcal{B} \;:=\; \mathbb{H}\otimes_{\mathbb{R}}\mathbb{C},
\end{equation}
where $\mathbb{H}$ is the real quaternion algebra and $\mathcal{B}$ is the
biquaternion algebra, isomorphic to $\mathrm{Mat}(2,\mathbb{C})$.
\end{definition}

The scalar phase right action on $\mathcal{B}$ defines a U(1) symmetry:
\begin{equation}\label{eq:U1action}
  \Theta(x) \;\mapsto\; \Theta(x)\,e^{i\alpha(x)}, \qquad \alpha(x)\in\mathbb{R}.
\end{equation}
This is an exact algebraic identity: the right multiplication $e^{i\alpha}$
lies in the centre $Z(\mathcal{B}) \cong \mathbb{C}$, so it commutes with all
gauge and algebraic operations on $\Theta$.

\begin{proposition}[U(1) is exact]
\label{prop:U1exact}
The transformation \eqref{eq:U1action} is an exact symmetry of the kinetic
term $\mathrm{Tr}[(\partial_\mu\Theta)^\dagger(\partial^\mu\Theta)]$:
\begin{equation}
  \mathrm{Tr}[(\partial_\mu(\Theta e^{i\alpha}))^\dagger
              (\partial^\mu(\Theta e^{i\alpha}))]
  = \mathrm{Tr}[(\partial_\mu\Theta)^\dagger(\partial^\mu\Theta)],
\end{equation}
for $\alpha = \text{const}$.
\end{proposition}

\begin{proof}
Since $e^{i\alpha}$ is a scalar and unitary, $(e^{i\alpha})^\dagger = e^{-i\alpha}$,
so $e^{-i\alpha}e^{i\alpha} = 1$.  Cyclicity of the trace then gives the result.
\end{proof}

\noindent For a \emph{local} function $\alpha(x)$, the kinetic term acquires
an extra term $\partial_\mu\alpha$ which must be compensated by a gauge field.
This is the standard Weyl gauge-principle argument, which introduces the
electromagnetic potential:
\begin{equation}\label{eq:Amu}
  A_\mu(x) \;:=\; \frac{1}{e}\,\partial_\mu\varphi(x),
\end{equation}
where $\varphi(x)$ is the phase of $\Theta$ and $e$ is the elementary charge.
Under $\varphi \to \varphi + e\,\alpha(x)$:
\begin{equation}
  A_\mu \;\to\; A_\mu + \partial_\mu\alpha(x).
\end{equation}
This is the standard gauge transformation of the electromagnetic four-potential.
The gauge-covariant derivative that renders the Lagrangian locally invariant is:
\begin{equation}\label{eq:covderiv}
  D_\mu\Theta \;:=\; \partial_\mu\Theta + ie\,A_\mu\,\Theta.
\end{equation}

% ------------------------------------------------------------------
\section{Step 2: Electromagnetic Field Tensor}
% ------------------------------------------------------------------

\noindent\textbf{Status: \PROVED{} \Lzero}

\noindent\textbf{Primary sources:}
\texttt{canonical/interactions/qed.tex}, \S1;
\texttt{canonical/bridges/Maxwell\_limit\_bridge.tex}, \S2.

\medskip

\begin{definition}[Field-strength tensor]
The electromagnetic field-strength tensor is the curvature of the U(1)
connection $A_\mu$:
\begin{equation}\label{eq:Fmunu}
  F_{\mu\nu} \;:=\; \partial_\mu A_\nu - \partial_\nu A_\mu.
\end{equation}
\end{definition}

\begin{proposition}[Gauge invariance of $F_{\mu\nu}$]
Under $A_\mu \to A_\mu + \partial_\mu\alpha$, the tensor $F_{\mu\nu}$ is invariant:
$F_{\mu\nu} \to F_{\mu\nu}$.
\end{proposition}
\begin{proof}
$(\partial_\mu A_\nu + \partial_\mu\partial_\nu\alpha)
- (\partial_\nu A_\mu + \partial_\nu\partial_\mu\alpha)
= F_{\mu\nu}$, since partial derivatives commute.
\end{proof}

In the biquaternionic formalism, the same object appears as the imaginary part
of the field curvature:
\begin{equation}
  \mathcal{F} \;=\; (\mathbf{E} + i\,\mathbf{B})\cdot\boldsymbol{\sigma},
\end{equation}
where $\mathbf{E}$ and $\mathbf{B}$ are the electric and magnetic field vectors
and $\boldsymbol{\sigma}$ are the Pauli matrices
(\texttt{consolidation\_project/appendix\_C\_electromagnetism\_gauge\_consolidated.tex}).
The standard tensor components are:
\begin{equation}
  F_{0i} = E_i, \qquad F_{ij} = -\varepsilon_{ijk}\,B_k.
\end{equation}

\begin{proposition}[Bianchi identity]
\label{prop:bianchi}
$F_{\mu\nu}$ satisfies:
\begin{equation}\label{eq:bianchi}
  \partial_\lambda F_{\mu\nu} + \partial_\mu F_{\nu\lambda}
  + \partial_\nu F_{\lambda\mu} = 0.
\end{equation}
\end{proposition}
\begin{proof}
Direct substitution of \eqref{eq:Fmunu} and use of symmetry of second
partial derivatives.
\end{proof}

% ------------------------------------------------------------------
\section{Step 3: Maxwell Equations}
% ------------------------------------------------------------------

\noindent\textbf{Status: \PROVED{} \Lone}

\noindent\textbf{Primary sources:}
\texttt{canonical/interactions/qed.tex}, \S3;
\texttt{canonical/bridges/Maxwell\_limit\_bridge.tex}, \S3--\S4.

\medskip

The QED Lagrangian density in the UBT framework is:
\begin{equation}\label{eq:LQED}
  \mathcal{L}_{\mathrm{QED}}
  \;=\;
  \mathrm{Tr}\bigl[(D_\mu\Theta)^\dagger(D^\mu\Theta)\bigr]
  \;-\;
  \frac{1}{4}\,F_{\mu\nu}F^{\mu\nu}.
\end{equation}

\subsection{Homogeneous Maxwell equations (Bianchi)}

Proposition~\ref{prop:bianchi} already gives the homogeneous equations.
Translating to three-vector notation:
\begin{align}
  \nabla\cdot\mathbf{B} &= 0, \label{eq:max1} \\
  \nabla\times\mathbf{E} + \frac{\partial\mathbf{B}}{\partial t} &= 0. \label{eq:max2}
\end{align}
These are \emph{automatic consequences} of defining $F_{\mu\nu}$ as a
curl (Definition~\ref{eq:Fmunu}) and require no dynamics. Status: \Lzero.

\subsection{Inhomogeneous Maxwell equations (Euler--Lagrange)}

\begin{theorem}[Inhomogeneous Maxwell equations]
\label{thm:maxwell_inhom}
The Euler--Lagrange equation of $\mathcal{L}_{\mathrm{QED}}$
with respect to $A_\nu$ is:
\begin{equation}\label{eq:maxwell_inhom}
  \partial_\mu F^{\mu\nu} \;=\; j^\nu,
\end{equation}
where $j^\nu$ is the conserved electromagnetic four-current.
\end{theorem}

\begin{proof}
The photon sector of \eqref{eq:LQED} is $\mathcal{L}_{\mathrm{ph}} = -\tfrac{1}{4}F_{\mu\nu}F^{\mu\nu}$.
We have:
\begin{equation}
  \frac{\partial\mathcal{L}_{\mathrm{ph}}}{\partial(\partial_\mu A_\nu)}
  = -F^{\mu\nu}.
\end{equation}
The Euler--Lagrange equation $\partial_\mu(\partial\mathcal{L}/\partial(\partial_\mu A_\nu)) = \partial\mathcal{L}/\partial A_\nu$
gives:
\begin{equation}
  -\partial_\mu F^{\mu\nu} = \frac{\partial\mathcal{L}_{\mathrm{matter}}}{\partial A_\nu}
  \;=:\; -j^\nu,
\end{equation}
yielding \eqref{eq:maxwell_inhom}. \qed
\end{proof}

\noindent In three-vector notation:
\begin{align}
  \nabla\cdot\mathbf{E} &= \rho, \label{eq:max3} \\
  \nabla\times\mathbf{B} - \frac{\partial\mathbf{E}}{\partial t} &= \mathbf{j}. \label{eq:max4}
\end{align}
Equations \eqref{eq:max1}--\eqref{eq:max4} are Maxwell's equations in
Gaussian units ($c=1$).

% ------------------------------------------------------------------
\section{Step 4: Fermion Coupling}
% ------------------------------------------------------------------

\noindent\textbf{Status: \PROVED{} \Lone}

\noindent\textbf{Primary sources:}
\texttt{canonical/interactions/qed.tex}, \S2 and \S4;
\texttt{unified\_biquaternion\_theory/ubt\_appendix\_10\_qed\_dirac.tex}.

\medskip

The spinor excitations of $\Theta$ decompose as:
\begin{equation}\label{eq:decomp}
  \Theta(q,\tau) \;=\; \psi_+(q,\tau) \otimes u_+ + \psi_-(q,\tau) \otimes u_-,
\end{equation}
where $u_\pm$ are the two-component spinor bases of $\mathrm{Mat}(2,\mathbb{C})$
and $\psi_\pm$ are four-component Dirac spinors.

The gauge-covariant derivative \eqref{eq:covderiv} acts on $\psi$ as:
\begin{equation}\label{eq:Dmu_psi}
  D_\mu\psi \;=\; \partial_\mu\psi - ie\,A_\mu\,\psi.
\end{equation}
This is the \emph{minimal coupling prescription}: the only coupling between
$\psi$ and $A_\mu$ consistent with local U(1) invariance and renormalizability.

\begin{proposition}[Gauge invariance of fermion kinetic term]
The fermion kinetic term $\bar\psi\,i\gamma^\mu D_\mu\psi$ is invariant under
simultaneous transformations $\psi\to e^{ie\alpha}\psi$ and
$A_\mu \to A_\mu + \partial_\mu\alpha$.
\end{proposition}
\begin{proof}
$\bar\psi\,i\gamma^\mu D_\mu(e^{ie\alpha}\psi)
= \bar\psi e^{-ie\alpha}\,i\gamma^\mu(e^{ie\alpha}\partial_\mu\psi
  + ie(\partial_\mu\alpha)e^{ie\alpha}\psi
  - ie(A_\mu+\partial_\mu\alpha)e^{ie\alpha}\psi)
= \bar\psi\,i\gamma^\mu(\partial_\mu\psi - ieA_\mu\psi)
= \bar\psi\,i\gamma^\mu D_\mu\psi$.
\end{proof}

The conserved Noether current is:
\begin{equation}\label{eq:jmu}
  j^\mu \;=\; e\,\bar\psi\,\gamma^\mu\psi,
\end{equation}
which is the standard Dirac current, satisfying $\partial_\mu j^\mu = 0$.

The complete QED Lagrangian including fermions is:
\begin{equation}\label{eq:LQED_full}
  \mathcal{L}_{\mathrm{QED}}
  \;=\;
  \bar\psi\,(i\gamma^\mu D_\mu - m)\,\psi
  \;-\;
  \frac{1}{4}\,F_{\mu\nu}F^{\mu\nu},
\end{equation}
where $m$ is the fermion mass.

% ------------------------------------------------------------------
\section{Step 5: Dirac Equation Limit}
% ------------------------------------------------------------------

\noindent\textbf{Status: \PROVED{} \Lone}

\noindent\textbf{Primary sources:}
\texttt{canonical/interactions/qed.tex}, \S4;
\texttt{unified\_biquaternion\_theory/ubt\_appendix\_10\_qed\_dirac.tex}, \S3.

\medskip

\begin{theorem}[Dirac equation from UBT]
\label{thm:dirac}
In the U(1) gauge sector and at constant phase ($\partial_\psi = 0$),
the Euler--Lagrange equation of \eqref{eq:LQED_full} with respect to
$\bar\psi$ is:
\begin{equation}\label{eq:dirac}
  (i\gamma^\mu D_\mu - m)\,\psi \;=\; 0,
  \qquad
  D_\mu = \partial_\mu - ie\,A_\mu.
\end{equation}
\end{theorem}

\begin{proof}
Varying the full action $S = \int d^4x\,\mathcal{L}_{\mathrm{QED}}$ with
respect to $\bar\psi$ gives:
\begin{equation}
  \frac{\delta S}{\delta\bar\psi}
  = (i\gamma^\mu D_\mu - m)\,\psi = 0,
\end{equation}
which is the Dirac equation \eqref{eq:dirac}. \qed
\end{proof}

\begin{remark}[Origin from UBT kinetic term]
In the UBT notation, the Dirac equation arises from the master equation
$\nabla^\dagger\nabla\Theta = \kappa\mathcal{T}$ in the U(1) sector:
\begin{equation}
  \nabla^\dagger\nabla\Theta \xrightarrow[\text{U(1) limit}]{\psi\to 0}
  (i\gamma^\mu D_\mu - m)\psi = 0.
\end{equation}
The imaginary time component $\psi = 0$ extracts the standard Dirac equation
from the biquaternionic master equation. The full biquaternionic equation
contains additional phase-curvature terms which vanish in this limit.
\end{remark}

\begin{remark}[Massive vs.\ massless fermion]
The mass term $m$ in \eqref{eq:dirac} is a free parameter in the classical
Dirac equation. In UBT the mass arises from the potential $V(\Theta)$ and is
related to the imaginary curvature of the $\Theta$ field. The precise derivation
of fermion masses is documented in \texttt{research/fermion\_mass\_program.md}
and is currently semi-empirical (see also \texttt{research/fermion\_mass\_status.md}).
\end{remark}

% ------------------------------------------------------------------
\section{Connection to Maxwell's Equations (Classical Limit)}
% ------------------------------------------------------------------

\noindent\textbf{Status: \PROVED{} \Lone}

In the classical (non-quantum) limit, one sets $\psi = 0$ (no fermions) in
\eqref{eq:LQED_full}, leaving only the Maxwell sector:
\begin{equation}
  \mathcal{L}_{\mathrm{Maxwell}} = -\frac{1}{4}F_{\mu\nu}F^{\mu\nu}.
\end{equation}
The Euler--Lagrange equation gives
$\partial_\mu F^{\mu\nu} = 0$ (vacuum Maxwell equations)
or $\partial_\mu F^{\mu\nu} = j^\nu$ with a classical current source.
This limit is proved in detail in \texttt{canonical/bridges/Maxwell\_complete\_derivation.tex}.

% ------------------------------------------------------------------
\section{Complete Chain Summary}
% ------------------------------------------------------------------

\begin{center}
\begin{tabular}{llll}
\hline
\textbf{Step} & \textbf{Input} & \textbf{Output} & \textbf{Status} \\
\hline
1 & $\Theta \in \mathcal{B}$, right-phase action & U(1) gauge symmetry & \PROVED{} \Lzero \\
2 & U(1) connection $A_\mu$ & $F_{\mu\nu}$, Bianchi identity & \PROVED{} \Lzero \\
3 & $\mathcal{L}_{\mathrm{QED}}$, var.\ wrt $A_\nu$ & Maxwell equations & \PROVED{} \Lone \\
4 & Spinor sector of $\Theta$, $D_\mu$ & Minimal fermion coupling & \PROVED{} \Lone \\
5 & Var.\ of $\mathcal{L}_{\mathrm{QED}}$ wrt $\bar\psi$ & Dirac equation & \PROVED{} \Lone \\
\hline
\end{tabular}
\end{center}

All five steps of the QED chain are established at proof level \Lone{} or
better. The chain is complete and self-contained: no step relies on an open
problem.

% ------------------------------------------------------------------
\section{Cross-References}
% ------------------------------------------------------------------

\begin{itemize}
  \item \textbf{U(1) gauge group canonical derivation}:
        \texttt{canonical/interactions/sm\_gauge.tex}, Theorem G.3
  \item \textbf{QED Lagrangian (canonical definition)}:
        \texttt{canonical/interactions/qed.tex}
  \item \textbf{Maxwell limit (detailed derivation)}:
        \texttt{canonical/bridges/Maxwell\_complete\_derivation.tex}
  \item \textbf{Maxwell limit (navigation bridge)}:
        \texttt{canonical/bridges/Maxwell\_limit\_bridge.tex}
  \item \textbf{QED limit status and open problems}:
        \texttt{canonical/bridges/QED\_limit\_bridge.tex}
  \item \textbf{Gauge group emergence}:
        \texttt{canonical/bridges/gauge\_emergence\_bridge.tex}
  \item \textbf{GR recovery chain (analogous)}:
        \texttt{canonical/bridges/GR\_chain\_bridge.tex}
  \item \textbf{Dirac equation (full treatment)}:
        \texttt{unified\_biquaternion\_theory/ubt\_appendix\_10\_qed\_dirac.tex}
\end{itemize}

\end{document}
